% ===============================================================
% tensor — a canonical-reference C++20 named-axis tensor library
% Workshop paper draft (target: NeurIPS / ICLR ML systems workshop)
% Snapshot: 2026-05-14
%
% Status: DRAFT (Phase 7a M3 deliverable, A-stage of A→B staging).
% Empirical sections marked TODO until Phase 7a M2 micro-bench
% data lands.  Template assumes generic article.cls; swap for the
% target venue's style file before submission.
% ===============================================================

\documentclass[11pt,a4paper]{article}

\usepackage[margin=1in]{geometry}
\usepackage{lmodern}      % scalable fonts so microtype works
\usepackage{microtype}
\usepackage{hyperref}
\usepackage{url}
\usepackage{booktabs}
\usepackage{listings}
\usepackage{xcolor}
\usepackage{amsmath,amssymb}
\usepackage{cite}

% Lightweight \todo macro (avoids the heavyweight todonotes package).
% Renders as a coloured inline marker; suitable for a draft.
\newcommand{\todo}[1]{\textcolor{red}{\textbf{[TODO:}\ #1\textbf{]}}}

\hypersetup{
  colorlinks=true,
  linkcolor=black,
  citecolor=black,
  urlcolor=blue,
}

\lstset{
  basicstyle=\ttfamily\footnotesize,
  keywordstyle=\bfseries,
  commentstyle=\itshape\color{gray!60!black},
  showstringspaces=false,
  language=C++,
  morekeywords={constexpr,concept,requires,consteval,co_await,co_yield,co_return},
  breaklines=true,
  frame=single,
  framesep=4pt,
  framerule=0.3pt,
  rulecolor=\color{gray!40},
  xleftmargin=4pt,
}

% =========================================================
\title{\texttt{tensor}: A C++20 Named-Axis Differentiable
  Tensor Library with a \LaTeX-Subset User-Defined Literal as
  First-Class Expression Surface}

\author{%
  uyuutosa\thanks{Solo maintainer; affiliation TBD. GitHub:
    \url{https://github.com/uyuutosa/tensor}}%
}

\date{2026-05-14 (working draft; not yet submitted)}

\begin{document}

\maketitle

\begin{abstract}
We present \texttt{tensor}, a C++20 header-only named-axis
differentiable tensor library, and argue for its position in the
2026 named-axis-tensor landscape as a specialist canonical
reference rather than a production-scale competitor. Three design
elements distinguish \texttt{tensor} from existing libraries:
(i) a \emph{hybrid named-axis API} that pairs compile-time
non-type template parameters (\texttt{LabelTag<"i">}) with a
runtime \texttt{DynamicShape}, giving callers both static
axis-name checking and ergonomic data-driven code from the same
type universe; (ii) a \LaTeX-subset user-defined literal
(\texttt{\_tex} UDL) that parses a fragment of Einstein-style
notation at C++ compile time, making the formula literally the
program; (iii) a Hexagonal-lite \texttt{KernelBackend} port with
three swappable adapters (reference, Eigen, WebGPU via Dawn) as
an explicit teaching surface for architecture-as-pedagogy. We
position these choices against eight contemporary
named-axis-adjacent libraries (xarray, PyTorch named tensors,
einops, einx, Haliax, Penzai, Nx, tinygrad), each of which
overlaps with at most three of \texttt{tensor}'s five
characteristic attributes. Production-scale ML training is an
explicit non-goal; we redirect that audience to JAX-side
libraries (Haliax in particular) and concentrate
solo-maintainer bandwidth on the moats that remain
uncontested. We discuss the design choices, sketch a worked
example from multi-view geometry, and outline open problems that
future implementers porting the substrate to alternative
languages (Rust, Zig, Mojo) would inherit.
\end{abstract}

% =========================================================
\section{Introduction}
\label{sec:intro}

Named-axis tensor algebra has accumulated steady literature over
the last five years. Chiang and Rush~\cite{ChiangRush2021}
formalised the notation as a mathematical language; Rogozhnikov's
einops~\cite{einops2022} demonstrated that a notation layer can
gain widespread adoption when applied across multiple host
frameworks; the einx successor~\cite{einx2026} extended this
substantially. In parallel, two named-axis libraries native to
the JAX ecosystem matured to production-scale: Haliax (now part
of the Marin foundation-model effort) has trained 70B-parameter
language models on TPU v4-2048 substrates, and Penzai from Google
DeepMind has integrated named axes with PyTree-based notebook
ergonomics. PyTorch's own \texttt{torch.named\_tensor} surface,
introduced in 2019, remains officially
``experimental and prototype'' as of PyTorch 2.12.

Within this landscape (see Section~\ref{sec:related-work}), the
combination \emph{C++20 native, first-class named axes,
compile-time axis-name checking, and a \LaTeX-subset
expression surface} is not occupied by any existing library. We
argue this gap is interesting not because it represents a
commercially competitive niche\,---\,it does not\,---\,but
because it constitutes a defensible \emph{canonical-reference}
position: an artifact future implementers can cite when porting
the named-axis substrate to alternative ecosystems, an artifact
educators can use to demonstrate modern C++ tensor algebra
without entanglement in a production framework's complexity, and
an artifact textbook authors can quote as a worked example of
``the formula is the program''.

We make three contributions:
\begin{enumerate}
\item A \LaTeX-subset user-defined literal that parses tensor
expressions at C++ compile time
(\S\ref{sec:tex-udl});
\item A hybrid named-axis API combining compile-time NTTPs and
runtime \texttt{DynamicShape}
(\S\ref{sec:hybrid-axes});
\item A Hexagonal-lite kernel-backend port architecture with
three swappable adapters as a deliberate teaching surface
(\S\ref{sec:kernel-port}).
\end{enumerate}

Section~\ref{sec:implementation} sketches the implementation;
Section~\ref{sec:case-study} works a multi-view-geometry case
study; Section~\ref{sec:related-work} positions
\texttt{tensor} against the eight contemporary libraries that
informed our design;
Section~\ref{sec:conclusion} discusses open problems.

% =========================================================
\section{Background}
\label{sec:background}

\paragraph{Named tensor notation.} A tensor in the Chiang--Rush
formulation~\cite{ChiangRush2021} carries each axis with an
explicit name, e.g.\ $A_{\textrm{batch},\textrm{seq},\textrm{embed}}$,
rather than a positional index. A contraction
$C_{ij} = \sum_k A_{ik} B_{kj}$ is then unambiguous regardless of
the order in which the operands' axes are laid out in memory. The
algebraic gain is the elimination of an entire class of
``transposition bugs'' that arise when positional axes are
re-shuffled silently by broadcasting rules. The implementation
gain is the ability to typecheck axis-name agreement at every
operation site.

\paragraph{Compile-time vs runtime axis checking.} Existing
named-axis libraries check axis-name agreement at one of three
points in time:

\begin{itemize}
\item \emph{At Python interpretation time}, when each operator
call inspects its inputs' name metadata (xarray, PyTorch named
tensors).
\item \emph{At JAX trace time}, when the first \texttt{jit}-ed
invocation compiles the computation graph (Haliax, Penzai).
\item \emph{Not at all}; named-axis surfaces are notation
conventions only (einops, einx, tinygrad).
\end{itemize}

C++20's non-type template parameters (NTTPs) admit a fourth
point: \emph{at C++ compile time}, before the program runs at
all. \texttt{tensor} uses this option for its compile-time-checked
path, falling back to the second option (a runtime check that
nonetheless happens once per operator, before any tensor
arithmetic) for its dynamic path.

% =========================================================
\section{The \texttt{\_tex} User-Defined Literal}
\label{sec:tex-udl}

The \texttt{\_tex} UDL parses a subset of \LaTeX{} tensor
notation at compile time into a typed AST. The motivating idiom is:

\begin{lstlisting}
auto expr = R"(c_{ij} = a_{ik} b_{kj})"_tex;
// expr is a parsed Expression; no runtime parser invocation.

auto C = tensor::tex::Evaluator<double>{}.eval(expr,
    {{"a", A}, {"b", B}});
// C is a DynamicTensor<double> with axes (i, j).
\end{lstlisting}

\noindent The user writes mathematics verbatim. The compiler
parses the literal into a \texttt{tex::Expression} via a
\texttt{consteval}-friendly recursive-descent parser; the
resulting AST is then handed to an \texttt{Evaluator} at runtime,
which routes each algebraic operation through the active
\texttt{KernelBackend} adapter
(\S\ref{sec:kernel-port}).

The supported subset as of 2026-05-14 includes assignment, sum
over repeated indices, scalar coefficients, parenthesisation, and
ordinary arithmetic operators. Phase 7c work~\cite{tensorRoadmap2026}
will extend the subset to product reductions ($\prod_i$),
Kronecker delta ($\delta_{ij}$), and the gradient operator
($\nabla_x$); the gradient extension will be of particular
interest because it composes with the tape-based reverse-mode
autograd subsystem~\cite{tensorADR0007} to make
$\nabla_x f(x)_{\dots}$ a compile-time-typed, autograd-routed
expression\,---\,a feature, to our knowledge, not implemented by
any other named-axis library.

\paragraph{Why the formula \emph{is} the program.} A persistent
gap between mathematical exposition and reference implementation
is that the implementation has to translate the math into a
programming-language idiom that obscures the original notation.
einops's dim-strings (\verb|"b h w c -> b c h w"|) make a
specific case of this translation explicit and compact, but
remain a notation distinct from the math literature's. The
\LaTeX-subset choice here is deliberate: it matches the
notation a researcher would write on a whiteboard, in a paper,
or in a textbook chapter without modification, eliminating the
mental round-trip.

% =========================================================
\section{Hybrid Compile-Time \& Runtime Named Axes}
\label{sec:hybrid-axes}

\texttt{tensor} exposes two coordinated axis-name surfaces:

\paragraph{Compile-time path.} A \texttt{LabelTag<"i">}
non-type template parameter carries an axis name as part of a
type. Contractions check name agreement at template
instantiation:

\begin{lstlisting}
template <auto I, auto J, auto K>
auto matmul(Tensor<double, Axis<I>, Axis<K>> a,
            Tensor<double, Axis<K>, Axis<J>> b)
    -> Tensor<double, Axis<I>, Axis<J>>;

// matmul(A_ij, B_jk) compiles.
// matmul(A_ij, B_kj) fails at template substitution.
\end{lstlisting}

This path is intended for fully-typed pipelines where the
axis structure is known by the call site.

\paragraph{Runtime path.} A \texttt{DynamicShape} carries axis
names and extents as runtime data:

\begin{lstlisting}
auto A = DynamicTensor<double>(
    DynamicShape{{"batch", 32}, {"seq", 128}, {"embed", 768}},
    /*data*/);

auto contracted = contract(A, B, {"seq"});
// Throws if "seq" is not present in both tensors.
\end{lstlisting}

This path supports data-driven code paths where the axis
structure depends on configuration. The two paths share the same
underlying memory representation and produce identical numerical
output for shared workloads, so a user can begin with the
ergonomic runtime path and refactor performance-critical inner
loops to the compile-time path without changing memory layout or
arithmetic semantics.

Other named-axis libraries~\cite{Haliax2024,Penzai2024,einx2026}
generally pick one side of this trade-off: typically runtime
checks, with the JAX-based ones gaining trace-time checks as a
secondary stage. The \texttt{tensor} hybrid surface is, to our
knowledge, the first in production-quality form.

% =========================================================
\section{Kernel-Backend Port as Teaching Surface}
\label{sec:kernel-port}

The core's tensor operations factor through a
\texttt{KernelBackend} port comprising fifteen
methods~\cite{tensorADR0011}: element-wise primitives,
contractions, reductions, broadcasts, and a few utility
operations. Three concrete adapters implement this port:

\begin{itemize}
\item \emph{Reference}: pure-C++, canonical answer; used as the
oracle in all cross-validation tests~\cite{tensorADR0009}.
\item \emph{Eigen}: an Eigen-3.4 adapter for SIMD-vectorised CPU
arithmetic.
\item \emph{WebGPU}: a Dawn-backed WGSL-kernel
adapter~\cite{tensorADR0006,tensorADR0016} for GPU acceleration
that ships from the same wheel on Linux, macOS,
and (planned) Windows.
\end{itemize}

\noindent This Hexagonal-lite~\cite{tensorADR0009} pattern is
not unusual on its own; the choice we wish to highlight is that
the port surface is itself the \emph{teaching artifact}. The
tutorial corpus that ships with the library~\cite{tensorTutorials}
walks a reader through swapping backends, observing that the
arithmetic semantics are invariant under backend selection, and
working out the boundary between the substrate-agnostic Domain
core and the substrate-specific adapter. JAX and PyTorch's
analogous abstractions (XLA, ATen, etc.) are similarly powerful
but largely opaque to users; \texttt{tensor}'s deliberately
simple port is intended as a worked example of the architectural
idea.

% =========================================================
\section{Implementation}
\label{sec:implementation}

The C++20 core is header-only~\cite{tensorADR0008} and depends
on \texttt{<mdspan>} for memory layout interop, vcpkg for
substrate management~\cite{tensorADR0003}, and (optionally) Dawn
for the WebGPU adapter. The Python SDK
binds via nanobind~\cite{tensorADR0018} and ships as a
PEP-508-extras-flavoured wheel set
(\texttt{tensor-named-axis[reference,eigen,webgpu,all]}); a
Phase~6.6 follow-up~\cite{tensorADR0020} will split the binding
into a type-owner module and per-backend kernel modules to
permit true in-process backend switching. Tape-based
reverse-mode autograd is exposed both from the C++ core and from
the Python SDK; cross-validated parity against
\texttt{numpy.einsum} and a hand-coded Jacobian
holds within $10^{-12}$ tolerance for
\texttt{double} and $10^{-5}$ for \texttt{float} across all
adapters~\cite{tensorADR0017}.

Project governance is documented as twenty-one architecture
decision records~\cite{tensorADRs} in MADR
v3.0~\cite{MADRv3} form within an arc42 §1--12
documentation system~\cite{arc42} alongside a Diátaxis-organised
user manual~\cite{Diataxis}. This documentation discipline is
itself one of the moats the library aspires
to~\cite{tensorADR0015,tensorADR0021}: when the substrate is
ported to another language, the ADR chain documents the design
reasoning that the code alone cannot capture.

% =========================================================
\section{Case Study: Bundle Adjustment with Named Axes}
\label{sec:case-study}

\todo{Section to be expanded with concrete numerical results
from the Phase~7a M2 benchmark pass. Below is the structural
outline; numbers will be drawn from
\texttt{python/notebooks/04\_python-bundle-adjustment-perspective.ipynb}
once the micro-benchmark harness against einops and Haliax
lands.}

Multi-view geometry's bundle adjustment optimises camera and
landmark parameters to minimise reprojection error across
multiple views of a scene. The reprojection residual has
natural named axes: $\textrm{view}$, $\textrm{landmark}$,
$\textrm{image\_coord}$, $\textrm{world\_coord}$. With
\texttt{tensor}'s \texttt{\_tex} UDL, the residual becomes:

\begin{lstlisting}
auto residual = R"(
  r_{v,l,c} = u_{v,l,c} - K_{c,d} R_{v,d,w} (X_{l,w} - t_{v,w})
)"_tex;
\end{lstlisting}

\noindent and the gradient with respect to camera pose
$(R, t)$ and landmark positions $X$ is obtained by routing the
expression through the autograd tape. Preliminary measurements
on a synthetic 50-view, 200-landmark scene show that the named-axis
expression compiles to within \todo{x.xx\%} of the runtime of an
equivalent hand-coded einsum, while catching one class of
axis-mislabelling error (here, accidentally contracting
$\textrm{view}$ with $\textrm{landmark}$) at compile time rather
than producing a silently-wrong result.

\todo{Add timing table with three rows (reference / Eigen /
WebGPU) and three columns
(\texttt{tensor} / einops on numpy / Haliax on JAX).}

% =========================================================
\section{Related Work}
\label{sec:related-work}

A companion landscape comparison study~\cite{tensorComparison2026}
surveys nine contemporary named-axis-adjacent libraries across
eleven evaluation axes. We summarise the most relevant comparisons.

\paragraph{Haliax.}~\cite{Haliax2024} A JAX-based first-class
named-axis library produced by Stanford CRFM that powers
production-scale LLM training (70B+ parameters on TPU v4-2048)
within the Marin monorepo. Haliax shares the first-class
named-axis goal with \texttt{tensor} but differs in host
language (JAX/Python vs.\ C++) and lacks a compile-time-checked
or DSL-expression surface. \texttt{tensor}'s explicit non-goal of
production-scale ML competition redirects that audience to
Haliax.

\paragraph{Penzai.}~\cite{Penzai2024} Google DeepMind's JAX-based
named-axes library, distinguished by PyTree-based notebook
ergonomics. Same architectural posture as Haliax with respect to
this paper's distinguishing features.

\paragraph{einops / einx.}~\cite{einops2022,einx2026} Both are
notation translation layers that compile string-DSL expressions
to host-framework calls on NumPy, PyTorch, JAX, TF, or MLX. They
are notation surfaces over positional tensor types rather than
named-axis tensor libraries proper. einx (ICLR 2026 oral)
materially expands the expressive range of einops. Both are
likely complementary to \texttt{tensor} rather than
substitutes; a future Python-side adapter exposing einops/einx
strings as a second authoring surface alongside \texttt{\_tex}
is plausible Phase~8+ work.

\paragraph{xarray.}~\cite{xarray2017} The oldest named-axis
library in Python (2014). Excellent
scientific-computing adoption (climate, geo).
Pure-Python, no autograd, no native GPU; complementary
to \texttt{tensor} but not directly comparable on architectural
attributes.

\paragraph{PyTorch named tensors.} A library-internal
named-axis surface inside PyTorch, ``experimental and prototype''
since 2019; treated here as a reference for ``what happens
when named axes are bolted onto a positional framework'' rather
than as a competitor.

\paragraph{tinygrad.}~\cite{tinygrad2024} Minimal positional-only
deep-learning library; cited here as a positioning reference for
educational-first framing without named axes. tinygrad
demonstrates that an explicit educational target can sustain a
community and reach production demos
(Llama~3 inference) even with a positional API.

\paragraph{What is uncontested.} Crossing the comparison study's
matrix, the cell ``C++-native + first-class named axes +
compile-time axis-name checking + \LaTeX-subset DSL +
citation-grade documentation discipline'' contains exactly one
library as of 2026-05-14. We do not expect this to remain true
indefinitely, but the moats compound for as long as the cell
stays unoccupied.

% =========================================================
\section{Conclusion and Future Work}
\label{sec:conclusion}

\texttt{tensor} occupies a deliberately narrow position in the
2026 named-axis-tensor landscape: it is the C++-native
canonical-reference choice rather than the production-scale
choice. The strategic narrowing is recorded in the project's
ADR-0021~\cite{tensorADR0021} as an explicit response to the
2026-05 landscape in which Haliax and Penzai substantially fill
the production-Python-ML named-axis wedge that a year earlier
appeared still open. The architectural distinctives we have
described\,---\,the \texttt{\_tex} UDL, the hybrid
compile-time/runtime axis API, and the
\texttt{KernelBackend} port as teaching surface\,---\,are
preserved investments under that narrowing.

\paragraph{Open problems.} Three Phase 7+ directions remain. The
Phase~6.6 type-owner-separation refactor will close the only
known real Python-SDK tech debt
(R-P6.5.5~\cite{tensorADR0020}). The Phase~7c
\texttt{\_tex} parser
expansion~\cite{tensorRoadmap2026} will add the
$\nabla$, $\prod$, and $\delta$ operators, the first of which
composes with the autograd tape to produce
compile-time-typed automatically-differentiated expressions\,---
\,a feature unique in the surveyed landscape. A subsequent
empirical evaluation pass will replace the \texttt{TODO} markers
in Section~\ref{sec:case-study} with measured comparisons against
einops on NumPy and Haliax on JAX. We invite contributions
from the named-tensor community along all three directions.

\paragraph{Reproducibility envelope.} All claims in this paper
fall within the build-and-test
reproducibility envelope of the
\texttt{tensor} repository at commit
\todo{commit hash on submission}, under the bibliography
discipline documented in
ADR-0017~\cite{tensorADR0017}. The Phase~7a
landscape comparison study~\cite{tensorComparison2026}
under which the related-work claims sit will be
re-validated at the 2026-11-11 half-yearly recheck.

\section*{Acknowledgements}

\todo{Acknowledgements after submission decisions.}

% =========================================================
\bibliographystyle{plain}
\bibliography{refs}

\end{document}
